% Source PDF: source/普通高考/2018/2018北京文.pdf, page 1, question 3.
% PDF embedded-image attempt: pdfimages -list found no image objects; the source
% flowchart is vector artwork. Grid evidence: tmp/redraw_flowchart_2018_beijing_liberal/
% original_pages/page-1_grid100.jpg (100px major, 50px fine grid); the comparison
% original side is a no-grid crop from the same page render.
% Elements: 开始/结束 rounded boxes; initialization and update rectangles;
% increment rectangle; k>=3 decision diamond; 输出 s parallelogram; directed
% connectors and the loop-back branch labelled 否, with the terminating branch 是.
% Relations: execution flows top-to-bottom from 开始 through initialization,
% update, increment and the decision; 否 loops from the diamond to the update;
% 是 proceeds to output and then 结束. The update is
% s=s+(-1)^k·1/(1+k), exactly as printed in the source.
% solid segments: all flowchart borders, vertical/horizontal connectors and
% arrowheads. dashed segments: none. Labels are centred in their nodes; branch
% labels sit beside the corresponding connector without touching the borders.
% Directed edges: 开始→初始化框；初始化框→更新框；更新框→自增框；自增框→判定框；
% 判定框(是)→输出框→结束；判定框(否)→右侧回路→更新框入口。两条分支仅在
% 判定框处分叉，回路箭头独立指向更新框入口，不复用终止分支的线段。
\begin{tikzpicture}[
  x=1cm,
  y=0.8cm,
  >=Stealth,
  line width=0.45pt,
  flowtext/.style={font=\normalsize,anchor=center,align=center,
    execute at begin node={\setlength{\parindent}{0pt}}},
  every node/.style={flowtext},
  flowline/.style={-{Stealth[length=2.1mm,width=1.35mm]},line width=0.45pt},
  branchlabel/.style={flowtext,inner sep=0pt},
  flowbox/.style={flowtext,draw,minimum height=0.58cm,inner xsep=4pt,inner ysep=2pt},
  startstop/.style={flowbox,rounded corners=0.22cm},
  decision/.style={flowbox,diamond,aspect=2.65,minimum width=2.7cm,minimum height=0.88cm,inner xsep=4pt,inner ysep=2pt},
  io/.style={flowbox,trapezium,trapezium left angle=75,trapezium right angle=105,minimum width=1.75cm,minimum height=0.62cm,inner xsep=4pt,inner ysep=2pt}
]
  % Main nodes follow the source's vertical execution order.
  \node[startstop,minimum width=1.18cm] (start) at (0,5.95) {开始};
  \node[flowbox,minimum width=2.45cm] (init) at (0,4.74) {$k=1,s=1$};
  \node[flowbox,minimum width=3.8cm,minimum height=0.98cm] (update) at (0,3.12)
    {$\displaystyle s=s+(-1)^k\cdot\frac{1}{1+k}$};
  \node[flowbox,minimum width=1.82cm] (increment) at (0,1.72) {$k=k+1$};
  \node[decision] (test) at (0,0.30) {$k\geq 3$};
  \node[io] (output) at (0,-1.15) {输出 $s$};
  \node[startstop,minimum width=1.18cm] (finish) at (0,-2.45) {结束};

  % Straight-through execution arrows.
  \draw[flowline] (start.south) -- (init.north);
  \draw (init.south) -- (0,4.05);
  \draw[flowline] (0,4.05) -- (update.north);
  \draw[flowline] (update.south) -- (increment.north);
  \draw[flowline] (increment.south) -- (test.north);
  \draw[flowline] (test.south) -- node[right=1.5pt] {是} (output.north);
  \draw[flowline] (output.south) -- (finish.north);

  % The negative branch returns to the shared update-entry trunk; the main
  % trunk's directed arrow marks the single continuation into the update box.
  \draw (test.east) -- (2.30,0.30);
  \draw (2.30,0.30) -- (2.30,4.05) -- (0,4.05);
  \node[branchlabel,anchor=east] at (2.15,0.48) {否};
\end{tikzpicture}
